{\small\textbf{Critiques (Zheng, 1:00) [22:00–23:00]}}\par\smallskip
Fourth, the open questions, because this is not a clean win.
\begin{itemize}\setlength\itemsep{0pt}
\item \textbf{Discretion needs credibility.} Without a mechanical rule, markets must trust DNB's judgement. That is the classic rules-versus-discretion trade-off: a rule is predictable but can be wrong, as we saw on slide 9; discretion is flexible but must be explained, every quarter, to stay credible. DNB's quarterly dashboard is its answer.
\item \textbf{Calibration.} Why two percent? ECB estimates suggest 1.1 to 1.8.
\item \textbf{It has never been tested:} we do not yet know whether banks will lend released capital or hoard it.
\item \textbf{Baseline.} Measured against 2021 rather than 2019, it \emph{was} an increase.
\item \textbf{Competing claims.} The ECB lists the Netherlands among countries that taxed banks' excess profits in 2022–23, partly the same profits that fund the buffer.
\end{itemize}
